\chapter{Introducción}

La planificación de horarios en cualquier ambito constituye una tarea compleja que implica la coordinación de múltiples variables. En el caso de los entornos universitarios, deben tenerse en cuenta variables como la disponibilidad del profesorado, la asignación de aulas y los grupos de estudiantes. Este proceso, cuando se realiza de forma manual, puede resultar ineficiente, propenso a errores y difícil de mantener ante cambios de última hora.

La generación de horarios presenta dificultades relacionadas con la gestión simultánea de múltiples restricciones y la necesidad de optimizar recursos limitados, lo que puede derivar en conflictos de solapamiento, infrautilización de espacios o distribución poco equilibrada de las clases.

\section{Justificación}

\section{Objetivo del Proyecto}

\section{Alcance del Proyecto}

\subsection{Estilo del Título 3}

Los estudiantes que lo deseen pueden hacer un uso libre de esta plantilla (respetando márgenes, encabezados y pies de página).

\section{Sobre los Proyectos Fin de Grado}

El Trabajo Fin de Grado (TFG) o su equivalente internacional, el \textit{Final Project}, es una materia con unas características particulares y de interés pedagógico para docentes e investigadores, por la complejidad que encarna. 

Estudios recientes apoyan a los estudiantes en su escritura académica mientras que otros se centran en la mejora de los resultados mediante la evaluación continua y uso intensivo del campus virtual y las estrategias de supervisión que emplean los docentes, entre otros aspectos.

Una diferencia esencial con los trabajos que se realizan en el resto de las asignaturas es que en el TFG se debe producir una integración del conocimiento, lo que supone al menos un peldaño de dificultad, especialmente en los sistemas educativos en los que predomina la fragmentación del currículum.

El TFG representa un salto cualitativo desde la división del saber, hacia el ejercicio profesional para el que normalmente se habilita con el título, que requiere del dominio de competencias que integran las diferentes parcelas del conocimiento teórico-práctico adquirido a lo largo de la carrera \cite{definition}.

\subsection{Objetivos}

El Proyecto Fin de Grado (en adelante PFG) es una de las actividades de aprendizaje más importantes del programa de los Grados en Ingeniería. Por una parte, el estudiante tiene la oportunidad de profundizar en el estudio de un tema de su interés. Por otra, le permite desarrollar competencias y habilidades fundamentales, tales como la capacidad de resolver problemas, defender propuestas mediante una comunicación eficiente, escribir una documentación, y aplicar métodos y técnicas aprendidas durante el Grado, entre otras.

El PFG consiste en la planificación, realización, presentación y defensa de un proyecto o trabajo de investigación sobre un área específica de las áreas de conocimiento del grado en ingeniería cursado. Su finalidad es propiciar la aplicación de las habilidades y los conocimientos adquiridos en el resto de la carrera.

Se intenta simular en lo posible la situación típica en que se va a ver envuelto el nuevo profesional al salir del entorno universitario, enlazando el mundo universitario con el laboral e integrando en un trabajo los diferentes conocimientos adquiridos.

El planteamiento del PFG debe ser, por tanto, práctico, lo más práctico posible, bajo el supuesto de que la teoría, así como las prácticas de cada asignatura, ya se conocen. Y no sólo práctico, sino también “comercial”, ya que ésta va a ser la práctica habitual de comportamiento una vez abandonado el mundo universitario y empezando a relacionarse en el mundo laboral.

El nuevo profesional tendrá que “vender” sus servicios a la organización a la que pertenezca (clientes internos) y a su vez los servicios y productos de esta organización a los clientes externos. El fruto de su trabajo no sólo debe tener calidad, sino que además debe saber comunicar con el entorno de forma satisfactoria.
Esta experiencia engloba dos facetas: la escrita y la oral.

Se trata, pues, por un lado, de que el alumno ensaye de la forma más realista posible la resolución de un caso práctico global y la elaboración de la documentación correspondiente de manera correcta, presentable y elegante, así como de que ensaye el modo de dirigirse ante un público y “comunicarse” con él, tarea también muy habitual y necesaria en el mundo empresarial.


\begin{quote}
  \textbf{El PFG debe considerarse no como un coste necesario para obtener el título, sino como una inversión que le servirá de modelo de entrenamiento para su ejercicio profesional.}
\end{quote}

Para terminar con los estilos de esta plantilla, también vamos a añadir los estilos para un nivel de anidamento de 4 niveles y de 5 niveles, por si algún estudiante los necesita.


\subsubsection{Estilo de Título 4}
Este es el estilo de título de nivel 4 por si es es necesario desarrollar algún punto con más detalle.

\paragraph{Estilo de título 5}
Este es el estilo de título de nivel 5 y es el último que se aporta en este ejemplo de plantilla de memoria. Si el estudiantado necesita desarrollar más niveles, pueden adaptar los estilos estándar que ya aparecen en el documento.

\chapter{Capítulo Segundo}
Si pensamos en las competencias genéricas que se deben trabajar en el PFG, llegamos a la siguiente enumeración:
\begin{enumerate}
  \item CG2. Comunicación verbal, Nivel 2: \dots
  \item CG3. Comunicación escrita, Nivel 2: \dots
  \item CG6: Sentido ético, Nivel 2: \dots
  \item CG8. Resolución de problemas, Nivel 2: \dots
\end{enumerate}

\section{Sobre Resultados de Aprendizaje}
\begin{itemize}
  \item CE1. Concebir y defender proyectos \dots
  \item CE2. Medición y evaluación técnica \dots
  \item CE3. Información y aprendizaje continuo \dots
\end{itemize}

\chapter{Valoración Ética del Proyecto}

Es importante resaltar que desde este curso, los estudiantes debéis añadir una valoración ética de vuestro proyecto, siguiendo las directrices que habéis recibido y trabajado en la asignatura de Ética.

